\documentclass{rapportECL}
\usepackage{lipsum}
\title{Note_de_lecture_TONGUE_KEVIN} %Titre du fichier
\begin{document}
% Désactivation de la page de garde et table des matières
\titre{Note de Lecture } %Titre du fichier .pdf
\UE{SHS-Management S9} %Nom de la UE
\sujet{Autonomie, apprentissage et responsabilité dans les organisations contemporaines} %Nom du sujet
\enseignant{Patrick \textsc{LAURENT}} %Nom de l'enseignant
\eleves{Kevin \textsc{TONGUE}} %Nom de l'élève
\UEU{ENISE-5GM}
\fairemarges
\fairepagedegarde
\tabledematieres

\section{Introduction}
Cette note de lecture combine les analyses de cinq textes relatifs à l'EC « SHS Management S9 ». Chaque section présente un résumé des idées principales suivies d'un commentaire argumenté, en lien avec les enseignements du cours. Les textes portent sur l'autonomie au travail, le rôle des leaders dans l'apprentissage, l'innovation managériale chez Haier, la construction de la culture organisationnelle et la responsabilité corporative à l'ère digitale.

\section{Note de Lecture : « L’autonomie ne se décrète pas, elle s’apprend ! Une approche des processus d’autonomisation des salariés à l’aune du concept de capabilité » de Benoît Grasser et Florent Noël (Revue Française de Gestion, n° 309, 2023)}

\subsection{Résumé des Idées et Positions Présentées dans le Texte}
Le texte de Grasser et Noël propose une analyse nuancée de l’autonomie au travail en mobilisant le concept de capabilité d’Amartya Sen, qui désigne la capacité réelle d’un individu à réaliser des fonctionnements (accomplissements) qu’il valorise, en tenant compte des ressources, des facteurs de conversion personnels, sociaux et environnementaux. Les auteurs évitent deux extrêmes caricaturaux : une vision déterministe du travail taylorien (où l’autonomie est absente) et une vision libérée (où elle est totale, comme dans les modèles holacratiques). Au lieu de cela, ils examinent, via une étude de cas, ce qui se passe lorsque, dans un environnement prescrit, le management encourage l’autonomie pour atteindre les objectifs organisationnels.

Les données montrent que l’autonomie accordée ne se traduit en marges de manœuvre réelles et en initiatives que si des facteurs de conversion (comme les compétences, le soutien managérial ou les normes collectives) permettent aux salariés de viser des réalisations valorisées à la fois pour eux-mêmes (émancipation personnelle) et pour l’organisation (efficacité collective). L’autonomie n’est pas une propriété décrétée par l’organisation ni une simple mise en œuvre par les salariés ; c’est un processus dynamique, marqué par des tensions (entre émancipation et contrôle, ou entre visions divergentes du "travail bien fait"). Ce processus s’appuie sur des boucles d’apprentissage : activation (mise en œuvre initiale), entretien (maintien via des pratiques récurrentes) et développement (évolution par l’innovation et la résolution de problèmes). Les auteurs concluent que l’autonomie s’apprend collectivement, nécessitant un accompagnement pour transformer les ressources en capabilités effectives.

\subsection{Commentaire Argumenté}

\subsubsection{Définition et Justification d’une Problématique}
La problématique retenue est la suivante : Dans quelle mesure l’autonomie au travail, souvent promue comme un levier d’efficacité organisationnelle, peut-elle être conçue non comme une prescription managériale statique, mais comme un processus d’apprentissage collectif qui intègre tensions et facteurs de conversion pour générer des réalisations valorisées ?

Cette problématique est justifiée par le texte lui-même, qui critique les approches binaires (déterminisme vs. libération) et met en lumière les limites d’une autonomie "décretée" sans accompagnement. Elle est pertinente car elle interroge la faisabilité réelle de l’autonomie dans des contextes organisationnels complexes, où les salariés doivent naviguer entre contraintes et initiatives. De plus, elle s’appuie sur les enseignements de l’EC "SHS Management S9", notamment les théories des organisations qui soulignent que la performance repose sur des interactions collectives (comme dans l’école des relations humaines, où Mayo démontre que la productivité dépend de facteurs psychosociaux et de groupes informels, plutôt que de prescriptions formelles). Ainsi, la problématique permet d’explorer comment l’autonomie, loin d’être un "one best way" (critiqué par les théories de contingence structurelle), s’adapte aux variables discriminantes comme la taille de l’organisation ou la technologie (Blau, Woodward).

\subsubsection{Définition et Justification d’un Plan d’Argumentation}
Le plan d’argumentation se structure en trois parties : (1) L’autonomie comme processus dynamique plutôt que décret statique ; (2) Le rôle des tensions et des facteurs de conversion dans la transformation des ressources en capabilités ; (3) Les boucles d’apprentissage comme mécanisme d’entretien et de développement de l’autonomie.

Ce plan est justifié car il suit logiquement la structure du texte : il commence par la critique des visions extrêmes (partie 1), approfondit les mécanismes internes (partie 2), et conclut sur les implications pratiques (partie 3). Il est lié à la problématique en démontrant comment l’autonomie s’apprend, et il intègre des arguments des autres cours pour enrichir l’analyse (par exemple, les théories économiques de l’organisation de Williamson, qui insistent sur les coûts de transaction et la régulation par l’organisation face à l’incertitude, ou les approches empiristes comme celles de Sloan, qui prônent une décentralisation pour une prise de décision plus efficace).

\subsubsection{Déroulement Discursif du Plan d’Argumentation}
\begin{enumerate}
\item \textbf{L’autonomie comme processus dynamique plutôt que décret statique}  
Le texte argue que l’autonomie n’est pas une propriété imposée par le management (comme dans un modèle taylorien revisité), mais un processus émergent qui s’apprend. Dans l’étude de cas, lorsque le management demande plus d’autonomie pour atteindre des objectifs, cela ne suffit pas sans un apprentissage actif. Cela résonne avec l’école classique des organisations (Taylor, Fayol), critiquée dans les cours pour son "one best way" rigide, qui ignore la dimension temporelle et historique des systèmes (comme défini dans l’introduction aux théories des organisations : une coordination continue pour des objectifs partagés). Au contraire, les auteurs montrent que l’autonomie émerge d’interactions, évitant les mythes du salarié "presse-bouton" ou "libéré". Cela s’aligne avec les théories de la contingence (Blau, Aldrich), où la taille et la différenciation structurelle influencent les modes d’organisation : dans une grande firme, l’autonomie doit s’adapter pour réduire les ratios d’encadrement, transformant ainsi les prescriptions en processus adaptatifs.

\item \textbf{Le rôle des tensions et des facteurs de conversion dans la transformation des ressources en capabilités}  
Selon Grasser et Noël, l’autonomie requiert des facteurs de conversion (personnels comme les compétences, sociaux comme les normes collectives, environnementaux comme le soutien managérial) pour transformer les ressources (autonomie accordée) en réalisations valorisées (marges de manœuvre et initiatives). Sans cela, les tensions émergent : entre émancipation individuelle et contrôle organisationnel, ou entre visions du "travail bien fait". Cela illustre comment l’autonomie n’est pas absolue, mais négociée. Des arguments des cours renforcent cela : dans les théories économiques (Coase, Williamson), l’organisation substitue au marché pour gérer l’incertitude et les coûts de transaction (imprévisibilité, vérification), où les facteurs de conversion agissent comme des "sous-systèmes opérationnels autonomes" pour lutter contre la désorganisation. De même, l’école des relations humaines (Mayo, Bavelas) met l’accent sur les relations inter-individuelles et les réseaux de communication (chaîne, cercle, étoile), où les tensions psychosociaux influencent la performance collective, transformant les ressources en capabilités via des interactions informelles.

\item \textbf{Les boucles d’apprentissage comme mécanisme d’entretien et de développement de l’autonomie}  
Le texte identifie des boucles d’apprentissage : activation (mise en œuvre initiale via initiatives), entretien (maintien par pratiques récurrentes) et développement (évolution par résolution de problèmes). Cela permet non seulement d’activer l’autonomie, mais de la faire évoluer face aux tensions. Cela s’appuie sur les approches empiristes (Sloan, Drucker), où la culture organisationnelle (valeurs, rites, normes) mobilise les membres via une identité partagée, et où la décentralisation (centres de profit autonomes) favorise l’apprentissage par exception et intégration. Dans les théories de la contingence (Woodward), les technologies de production (unitaires vs. continues) influencent les modes d’organisation, suggérant que les boucles d’apprentissage adaptent l’autonomie aux contextes, évitant un "one best way" pour des valeurs médianes performantes.
\end{enumerate}

\subsubsection{Conclusion Persuasive}
En conclusion, l’autonomie au travail, loin d’être un décret managérial, émerge comme un processus d’apprentissage collectif intégrant tensions et facteurs de conversion pour des réalisations valorisées, comme le démontre le texte de Grasser et Noël. Cette vision persuasive invite les managers à favoriser des boucles d’apprentissage, alignées sur les théories organisationnelles vues en cours (relations humaines pour les interactions, contingence pour l’adaptation, empiristes pour la culture). Ainsi, promouvoir l’autonomie n’est pas une mode, mais une stratégie durable pour une performance collective, incitant à repenser les pratiques managériales pour qu’elles s’apprennent plutôt que s’imposent.

\section{Note de Lecture : « Leaders’ Critical Role in Building a Learning Culture » by Henrik Saabye and Thomas Borup (MIT Sloan Management Review, Spring 2025)}

\subsection{Résumé des Idées et Positions Présentées dans le Texte}
Le texte de Saabye et Borup met en lumière le rôle essentiel des leaders dans la construction d'une culture d'apprentissage organisationnel, vitale pour l'adaptation et la survie face aux disruptions technologiques, environnementales et aux demandes changeantes des clients. Les auteurs critiquent la délégation courante de l'apprentissage aux spécialistes en développement, qui se fait souvent hors contexte de travail (ateliers, formations), sans bâtir une capacité durable de résolution de problèmes. À travers deux études longitudinales chez Lego (où une équipe de design 3D fait face à un backlog croissant) et Velux, ils montrent que les leaders doivent adopter une approche contre-intuitive : "aller lentement pour aller vite", en priorisant le développement des compétences de résolution systématique de problèmes chez les employés.

Les leaders deviennent des facilitateurs d'apprentissage en encadrant les problèmes comme des opportunités de croissance, en enseignant des méthodes structurées comme l'A3 thinking (analyse systématique, causes racines, solutions). Cela crée un effet d'entraînement : les employés gagnent en confiance et en indépendance, résolvent les problèmes plus efficacement, tandis que les leaders affinent leurs compétences en coaching. Le texte illustre cela chez Lego, où les leaders ont transformé une équipe surchargée en une unité plus agile via un focus sur l'apprentissage contextuel. Chez Velux, des pratiques similaires ont renforcé l'efficacité organisationnelle. Les auteurs concluent que les leaders doivent investir du temps dans l'enseignement et le coaching pour une transformation durable, surmontant les défis comme la pression temporelle ou la résistance au changement.

\subsection{Commentaire Argumenté}

\subsubsection{Définition et Justification d’une Problématique}
La problématique retenue est la suivante : Dans quelle mesure les leaders, souvent perçus comme des décideurs stratégiques, doivent-ils assumer un rôle actif de facilitateurs d'apprentissage pour transformer les organisations en entités adaptatives, en intégrant des approches contre-intuitives comme "aller lentement pour aller vite" afin de développer une culture de résolution de problèmes durable ?

Cette problématique est justifiée par le texte, qui souligne les limites de la délégation externe de l'apprentissage et insiste sur le rôle central des leaders dans la facilitation contextuelle pour une adaptation organisationnelle. Elle est pertinente car elle interroge la transition des leaders d'un mode directif à un mode coachant, dans des contextes de haute volatilité. Elle s’appuie sur les enseignements de l’EC "SHS Management S9", notamment les théories empiristes (Sloan, Drucker), où la décentralisation et la gestion par exception favorisent une culture mobilisatrice (valeurs, rites), et les théories de contingence structurelle (Blau, Woodward), qui montrent que la taille et la technologie influencent les modes d'organisation, nécessitant un apprentissage adaptatif pour une performance médiane optimale.

\subsubsection{Définition et Justification d’un Plan d’Argumentation}
Le plan d’argumentation se structure en trois parties : (1) Le rôle contre-intuitif des leaders comme facilitateurs d'apprentissage plutôt que délégateurs ; (2) L'intégration des défis et méthodes structurées pour transformer les problèmes en opportunités de croissance ; (3) Les mécanismes d'effet d'entraînement et d'affinage des compétences pour une culture d'apprentissage durable.

Ce plan est justifié car il suit la logique du texte : critique initiale de la délégation (partie 1), approfondissement des pratiques chez Lego et Velux (partie 2), et implications à long terme (partie 3). Il est lié à la problématique en démontrant comment les leaders construisent une culture adaptative, et intègre des arguments des autres cours pour enrichir l'analyse (par exemple, les théories économiques de Williamson, qui soulignent la régulation organisationnelle face à l'incertitude via des sous-systèmes autonomes, ou l'école des relations humaines de Mayo, où les facteurs psychosociaux et les communications influencent la performance collective).

\subsubsection{Déroulement Discursif du Plan d’Argumentation}
\begin{enumerate}
\item \textbf{Le rôle contre-intuitif des leaders comme facilitateurs d'apprentissage plutôt que délégateurs}  
Le texte argue que les leaders doivent prioriser l'apprentissage contextuel plutôt que de le déléguer à des spécialistes externes, adoptant "aller lentement pour aller vite" pour bâtir une capacité durable. Chez Lego, cela a permis de réduire les backlogs en enseignant la résolution de problèmes in situ. Cela contredit l'école classique (Taylor, Fayol), vue en cours pour son focus sur la prescription formelle et le "one best way", ignorant les dimensions psychosociaux. Au contraire, les auteurs alignent sur les théories de la contingence (Blau, Aldrich), où la croissance organisationnelle accroît la différenciation structurelle, nécessitant des leaders facilitateurs pour maintenir un ratio d'encadrement efficace et adapter l'apprentissage à la taille et à la complexité.

\item \textbf{L'intégration des défis et méthodes structurées pour transformer les problèmes en opportunités de croissance}  
Selon Saabye et Borup, les leaders affrontent des défis comme la pression temporelle, mais en utilisant des méthodes comme l'A3 thinking (analyse des causes racines), ils transforment les problèmes en croissance. Cela requiert un coaching actif pour développer la confiance des employés. Des arguments des cours renforcent cela : dans les théories empiristes (Sloan), la décentralisation via centres de profit autonomes et évaluation par comparaison favorise une culture où les leaders arbitrent les conflits et gèrent par exception, intégrant les défis pour une identité organisationnelle forte (Schein : mythes, rites, valeurs). De même, l'école des relations humaines (Bavelas) met l'accent sur les réseaux de communication (indice de centralité, connexité), où les leaders facilitent les interactions pour surmonter les résistances, transformant les défis psychosociaux en opportunités collectives.

\item \textbf{Les mécanismes d'effet d'entraînement et d'affinage des compétences pour une culture d'apprentissage durable}  
Le texte identifie un effet ripple : les employés deviennent indépendants, et les leaders affinent leur coaching, créant une culture adaptable. Chez Velux, cela a renforcé l'efficacité via un focus sur l'apprentissage récurrent. Cela s'appuie sur les théories économiques (Coase, Williamson), où l'organisation substitue au marché pour gérer les coûts de transaction (imprévisibilité, vérification), et les boucles d'apprentissage agissent comme des sous-systèmes pour lutter contre la désorganisation environnementale. Dans les théories de contingence (Woodward), les technologies (production en série vs. continue) influencent les modes, suggérant que l'effet d'entraînement adapte la culture aux contextes, favorisant des valeurs médianes performantes via un affinage continu des compétences.
\end{enumerate}

\subsubsection{Conclusion Persuasive}
En conclusion, les leaders doivent embrasser un rôle actif de facilitateurs d'apprentissage pour bâtir une culture adaptative, comme le démontrent Saabye et Borup, en intégrant des approches contre-intuitives pour un effet durable. Cette vision persuasive appelle à un shift managérial, aligné sur les théories organisationnelles des cours (empiristes pour la mobilisation culturelle, contingence pour l'adaptation contextuelle, relations humaines pour les interactions). Ainsi, investir dans l'apprentissage n'est pas un coût, mais une stratégie essentielle pour la résilience organisationnelle, incitant les leaders à prioriser le coaching pour des transformations pérennes plutôt que des solutions rapides.

\section{Note de Lecture : « Management Innovation Made in China: Haier’s Rendanheyi » by Jedrzej George Frynas, Michael J. Mol, and Kamel Mellahi (California Management Review, Vol. 61, 2018)}

\subsection{Résumé des Idées et Positions Présentées dans le Texte}
Le texte de Frynas, Mol et Mellahi explore comment les entreprises des marchés émergents innovent en management pour naviguer dans un environnement VUCA (volatility, uncertainty, complexity, ambiguity). À travers le cas de Haier, une entreprise chinoise de fabrication d'appareils électroménagers, les auteurs montrent comment elle a développé la plateforme Rendanheyi pour passer d'une structure hiérarchique traditionnelle à une organisation entrepreneuriale en ligne, hautement réactive. Rendanheyi intègre des pratiques comme des micro-entreprises autonomes (ZZJYTs), une focalisation sur les besoins des utilisateurs, et une structure en réseau pour réduire la bureaucratie et favoriser l'innovation.

Les auteurs soulignent que le développement de Rendanheyi dépend de contextes organisationnel (transformation interne), compétitif (concurrence globale), institutionnel (réformes chinoises) et technologique (internet et digitalisation). Cela permet à Haier de gérer le VUCA en encourageant l'entrepreneuriat interne et l'adaptation rapide. Ils proposent un modèle étendu d'innovation managériale, où les managers peuvent appliquer ces pratiques pour innover dans des contextes similaires. Le texte conclut que les innovations comme Rendanheyi, issues des marchés émergents, offrent des leçons pour les firmes globales, en reliant structure, comportement et performance (paradigme SCP de Bain).

\subsection{Commentaire Argumenté}

\subsubsection{Définition et Justification d’une Problématique}
La problématique retenue est la suivante : Dans quelle mesure les innovations managériales comme Rendanheyi, développées dans des contextes émergents marqués par le VUCA, peuvent-elles transformer les organisations hiérarchiques en entités entrepreneuriales adaptatives, en intégrant des facteurs contextuels pour une performance durable ?

Cette problématique est justifiée par le texte, qui met en avant le rôle des contextes dans la création de Rendanheyi et critique les approches occidentales statiques, en soulignant l'adaptabilité des marchés émergents. Elle est pertinente car elle interroge l'applicabilité globale de telles innovations dans un monde incertain. Elle s’appuie sur les enseignements de l’EC "SHS Management S9", notamment le paradigme "Structure – Comportement – Performance" (Bain), où la structure industrielle détermine les comportements et performances, et les théories de contingence structurelle (Blau, Woodward), qui montrent que les variables comme la taille ou la technologie influencent les modes d'organisation.

\subsubsection{Définition et Justification d’un Plan d’Argumentation}
Le plan d’argumentation se structure en trois parties : (1) Rendanheyi comme innovation contextuelle face au VUCA, plutôt qu'un modèle universel ; (2) L'intégration des facteurs organisationnels, compétitifs, institutionnels et technologiques pour transformer la structure hiérarchique ; (3) Le modèle étendu d'innovation managériale comme mécanisme pour une adaptation et performance durables.

Ce plan est justifié car il suit la structure du texte : description du contexte VUCA et de Rendanheyi (partie 1), analyse des facteurs contextuels (partie 2), et implications pratiques via un modèle étendu (partie 3). Il est lié à la problématique en démontrant comment les innovations émergentes favorisent l'adaptabilité, et intègre des arguments des autres cours pour enrichir l'analyse (par exemple, les théories économiques de Williamson sur les coûts de transaction et la régulation organisationnelle, ou les approches empiristes de Sloan pour la décentralisation et les centres de profit autonomes).

\subsubsection{Déroulement Discursif du Plan d’Argumentation}
\begin{enumerate}
\item \textbf{Rendanheyi comme innovation contextuelle face au VUCA, plutôt qu'un modèle universel}  
Le texte argue que Rendanheyi émerge des défis VUCA des marchés émergents, transformant Haier d'une firme hiérarchique en une plateforme entrepreneuriale via des micro-entreprises autonomes. Cela contredit le "one best way" de l'école classique (Taylor, Fayol), critiqué en cours pour ignorer les contextes, et aligne sur les théories de contingence (Blau, Aldrich), où la taille croissante accroît la différenciation structurelle à un taux décroissant, nécessitant des innovations adaptatives pour gérer l'incertitude. Chez Haier, Rendanheyi répond au VUCA en favorisant l'autonomie, similaire à comment Woodward lie les technologies de production (unitaires vs. massives) à des modes organisationnels performants aux valeurs médianes.

\item \textbf{L'intégration des facteurs organisationnels, compétitifs, institutionnels et technologiques pour transformer la structure hiérarchique}  
Selon Frynas et al., Rendanheyi intègre des contextes multiples : organisationnel (réduction de la bureaucratie), compétitif (focus sur les utilisateurs), institutionnel (réformes chinoises) et technologique (digitalisation). Cela permet une transformation en réseau réactif. Des arguments des cours renforcent cela : dans les théories économiques (Coase, Williamson), l'organisation substitue au marché pour minimiser les coûts de transaction (imprévisibilité, vérification), où les facteurs contextuels agissent comme des sous-systèmes autonomes pour contrer la désorganisation environnementale. De même, les approches empiristes (Sloan, Gélinier) prônent une décentralisation via centres de profit et gestion par exception, intégrant les facteurs pour une culture organisationnelle mobilisatrice (Schein : valeurs, mythes, rites), transformant ainsi la hiérarchie en structure adaptative.

\item \textbf{Le modèle étendu d'innovation managériale comme mécanisme pour une adaptation et performance durables}  
Le texte propose un modèle étendu où Rendanheyi lie structure, comportement et performance, offrant des leçons applicables. Cela s'appuie sur le paradigme SCP (Bain), vu en cours, où la concentration sectorielle (HHI, CRm) influence les comportements pour une performance accrue. Dans les théories de la contingence (Woodward), les typologies technologiques favorisent des innovations pour des performances médianes, suggérant que le modèle étendu adapte Rendanheyi aux contextes, favorisant une adaptation durable via l'entrepreneuriat interne. L'école des relations humaines (Mayo, Bavelas) ajoute que les réseaux de communication (indice de centralité) renforcent les comportements collectifs, rendant le modèle un mécanisme pour une performance psychosociologiquement soutenable.
\end{enumerate}

\subsubsection{Conclusion Persuasive}
En conclusion, les innovations comme Rendanheyi transforment les organisations dans des contextes VUCA en intégrant facteurs contextuels pour une adaptabilité durable, comme le démontrent Frynas, Mol et Mellahi. Cette vision persuasive appelle à adopter des modèles étendus, alignés sur les théories organisationnelles des cours (contingence pour l'adaptation, économiques pour la régulation, empiristes pour la décentralisation). Ainsi, les managers des marchés émergents offrent des leçons globales, incitant à repenser les pratiques pour une performance non hiérarchique mais entrepreneuriale, plutôt que rigide.

\section{Note de Lecture : « Building Culture From the Middle Out » by Spencer Harrison and Kristie Rogers (MIT Sloan Management Review, Spring 2024)}

\subsection{Résumé des Idées et Positions Présentées dans le Texte}
Le texte de Harrison et Rogers explore le rôle crucial des leaders intermédiaires dans la construction d'une culture organisationnelle saine et vibrante. Les auteurs commencent par noter que, bien que tout le monde se dise responsable de la culture (ensemble de valeurs partagées guidant le travail), peu savent la gérer concrètement au-delà de pratiques génériques comme une porte ouverte ou des revues de performance. La recherche montre que la culture prédit la haute performance, et les managers intermédiaires doivent passer de l'endossement (soutenir les normes existantes) à l'enrichissement (soutenir les expressions des valeurs dans les équipes).

Ils distinguent la "big-C culture" (valeurs officielles, documents canonisés, formations formelles) de la "small-c culture" (interactions quotidiennes vibrantes, comme des échanges renforçant les tâches). Basé sur une étude dans une Fortune 100 reconnue pour son environnement de travail, avec 120 entretiens, les auteurs identifient que les top 10\% des managers (en rétention et performance) lient big-C et small-c. Stratégies d'endossement : célébrer et préserver des aspects sélectionnés (ex. : égalité via inclusion en réunions) et apprendre des pairs (réseaux, panels d'erreurs). Stratégies d'enrichissement : innover culturellement (ex. : intégrer bien-être dans rapports de ventes) et empower les employés à innover (ex. : questions ice-breaker en réunions). Cela crée une culture adaptable, avec bénéfices comme l'engagement, la rétention et la diffusion d'innovations. Les auteurs concluent que la culture n'est pas uniforme mais évolue via des déviations saines, comme dans un écosystème.

\subsection{Commentaire Argumenté}

\subsubsection{Définition et Justification d’une Problématique}
La problématique retenue est la suivante : Dans quelle mesure les leaders intermédiaires, souvent relégués à un rôle passif d'endossement, peuvent-ils activement construire une culture organisationnelle en liant valeurs officielles (big-C) et interactions quotidiennes (small-c), pour favoriser une adaptabilité et une performance durables face aux changements ?

Cette problématique est justifiée par le texte, qui critique l'assomption que la culture est gérée par le top management ou HR, et met en lumière le gap entre endossement et enrichissement. Elle est pertinente car elle interroge le rôle des managers intermédiaires dans une culture vivante, surtout en périodes de croissance ou crise. Elle s’appuie sur les enseignements de l’EC "SHS Management S9", notamment les approches empiristes (Sloan, Drucker, Gélinier), où la culture (mythes, rites, valeurs) mobilise via identité et décentralisation, et les théories de contingence structurelle (Blau, Woodward), qui montrent que la taille et la technologie influencent les différenciations structurelles, nécessitant une culture adaptable pour une performance médiane.

\subsubsection{Définition et Justification d’un Plan d’Argumentation}
Le plan d’argumentation se structure en trois parties : (1) La distinction big-C/small-c comme cadre pour un rôle actif des managers intermédiaires plutôt que passif ; (2) L'intégration de stratégies d'endossement et d'enrichissement pour lier valeurs et interactions face aux défis ; (3) Les mécanismes d'innovation et d'empowerment pour une culture adaptable et performante.

Ce plan est justifié car il suit la logique du texte : introduction au gap et distinction (partie 1), stratégies détaillées (partie 2), et implications (partie 3). Il est lié à la problématique en démontrant comment les managers construisent la culture, et intègre des arguments des autres cours pour enrichir l'analyse (par exemple, les théories économiques de Williamson sur la régulation via sous-systèmes autonomes face à l'incertitude, ou l'école des relations humaines de Mayo pour les impacts psychosociaux des interactions).

\subsubsection{Déroulement Discursif du Plan d’Argumentation}
\begin{enumerate}
\item \textbf{La distinction big-C/small-c comme cadre pour un rôle actif des managers intermédiaires plutôt que passif}  
Le texte argue que les managers intermédiaires doivent enrichir la culture en liant big-C (valeurs officielles) et small-c (interactions quotidiennes), au-delà d'un endossement passif. Cela contredit l'école classique (Taylor, Fayol), vue en cours pour son "one best way" hiérarchique, ignorant les variantes culturelles. Au contraire, les auteurs alignent sur les théories de contingence (Blau, Aldrich), où la croissance accroît la différenciation à un taux décroissant, nécessitant des managers actifs pour gérer l'encadrement et adapter la culture à la complexité, comme chez Woodward où les technologies influencent les modes pour des performances médianes via des interactions vibrantes.

\item \textbf{L'intégration de stratégies d'endossement et d'enrichissement pour lier valeurs et interactions face aux défis}  
Selon Harrison et Rogers, l'endossement (célébration, apprentissage des pairs) et l'enrichissement (innovations) intègrent défis comme les changements (acquisitions, pandémies). Des arguments des cours renforcent cela : dans les approches empiristes (Sloan), la décentralisation via centres de profit et arbitrage des conflits enrichit la culture (Schein : identité via croyances, rites), liant big-C à small-c pour mobilisation. De même, l'école des relations humaines (Bavelas) met l'accent sur les réseaux (indice de centralité, connexité), où les managers facilitent les interactions pour surmonter les résistances, transformant les défis en liens culturels, comme dans les panels d'erreurs pour vulnérabilité partagée.

\item \textbf{Les mécanismes d'innovation et d'empowerment pour une culture adaptable et performante}  
Le texte identifie l'innovation (expérimentations) et l'empowerment (idées des employés) comme mécanismes pour une culture évolutive, avec bénéfices comme rétention et diffusion. Cela s'appuie sur les théories économiques (Coase, Williamson), où l'organisation gère les coûts de transaction (imprévisibilité) via sous-systèmes, et les innovations agissent comme des "field trips" pour adaptabilité. Dans les théories de contingence (Woodward), les typologies favorisent des cultures médianes performantes, suggérant que l'empowerment adapte small-c aux contextes, favorisant une performance via camaraderie et créativité, même en cas d'échecs.
\end{enumerate}

\subsubsection{Conclusion Persuasive}
En conclusion, les leaders intermédiaires construisent une culture adaptable en liant big-C et small-c via endossement et enrichissement, comme le démontrent Harrison et Rogers, favorisant engagement et performance. Cette vision persuasive appelle à un rôle actif, aligné sur les théories organisationnelles des cours (empiristes pour mobilisation, contingence pour adaptation, relations humaines pour interactions). Ainsi, la culture n'est pas un mythe passif mais un écosystème vivant, incitant les managers à innover pour une résilience organisationnelle durable plutôt que rigide.

\section{Note de Lecture : « Corporate Responsibility Meets the Digital Economy » by Leena Lankoski and N. Craig Smith (California Management Review, Vol. 67, 2025)}

\subsection{Résumé des Idées et Positions Présentées dans le Texte}
Le texte de Lankoski et Smith examine comment la transformation digitale bouleverse la responsabilité corporative (CR), en revisitant ses fondations face à un environnement "sismique" marqué par des gains sociétaux mais aussi des défis. Les auteurs notent que la digitalisation impacte des domaines entiers (économie, stratégie, éthique), et appliquent trois questions centrales à la CR : "responsible for what?" (responsable de quoi ?), "responsible toward whom?" (envers qui ?), et "who is responsible?" (qui est responsable ?). Pour "responsible for what?", la digitalisation élargit les enjeux CR à de nouveaux domaines comme la privacy, l'IA éthique, et les impacts environnementaux des data centers. Pour "responsible toward whom?", elle étend les stakeholders à des entités distantes ou algorithmiques (ex. : utilisateurs anonymes, futurs générations affectées par l'IA). Pour "who is responsible?", elle complique l'attribution en raison des chaînes de valeur complexes et des algorithmes autonomes.

Les auteurs illustrent ces impacts via des exemples comme les défis éthiques de l'IA et les régulations émergentes. Ils concluent par des implications : pour les managers (intégrer CR dans la digitalisation), les advocates CR (adapter les normes), les scholars (revisiter théories), et les policymakers (réguler pour équilibrer innovation et responsabilité). Le texte insiste sur une CR adaptée au digital, évitant les assumptions obsolètes pour aligner sur les développements contemporains.

\subsection{Commentaire Argumenté}

\subsubsection{Définition et Justification d’une Problématique}
La problématique retenue est la suivante : Dans quelle mesure la digitalisation, en transformant les environnements organisationnels, nécessite-t-elle une revisite des fondations de la responsabilité corporative (CR) via les questions de responsabilité (quoi, envers qui, qui), pour adapter les pratiques managériales à un monde VUCA tout en équilibrant innovation et éthique ?

Cette problématique est justifiée par le texte, qui souligne la nécessité de revisiter la CR face aux changements "révolutionnaires" du digital, évitant des assumptions périmées. Elle est pertinente car elle interroge l'intégration de la CR dans la digitalisation, surtout dans des contextes complexes. Elle s’appuie sur les enseignements de l’EC "SHS Management S9", notamment les théories économiques de l'organisation (Coase, Williamson), où la régulation organisationnelle gère l'incertitude et les coûts de transaction, et les approches empiristes (Sloan, Schein), où la culture intègre valeurs éthiques pour mobiliser.

\subsubsection{Définition et Justification d’un Plan d’Argumentation}
Le plan d’argumentation se structure en trois parties : (1) La revisite des questions CR comme réponse aux transformations digitales plutôt qu'aux modèles traditionnels ; (2) L'intégration des impacts sur les enjeux, stakeholders et attribution pour équilibrer éthique et innovation ; (3) Les implications pratiques pour managers, advocates, scholars et policymakers comme mécanismes d'adaptation durable.

Ce plan est justifié car il suit la structure du texte : cadre des questions (partie 1), analyse des impacts (partie 2), et implications (partie 3). Il est lié à la problématique en démontrant comment la CR s'adapte au digital, et intègre des arguments des autres cours pour enrichir l'analyse (par exemple, les théories de contingence structurelle de Blau pour la différenciation face à la complexité, ou l'école des relations humaines de Mayo pour les impacts psychosociaux sur les stakeholders).

\subsubsection{Déroulement Discursif du Plan d’Argumentation}
\begin{enumerate}
\item \textbf{La revisite des questions CR comme réponse aux transformations digitales plutôt qu'aux modèles traditionnels}  
Le texte argue que la digitalisation élargit la CR au-delà des modèles classiques, via les questions "quoi, envers qui, qui", pour aborder privacy, IA et impacts environnementaux. Cela contredit l'école classique (Taylor, Fayol), vue en cours pour son "one best way" rigide, ignorant les contextes digitaux. Au contraire, les auteurs alignent sur les théories de contingence (Woodward), où les technologies (comme la production continue digitale) influencent les modes organisationnels, nécessitant une revisite pour des performances médianes éthiques, adaptant la CR à la volatilité VUCA.

\item \textbf{L'intégration des impacts sur les enjeux, stakeholders et attribution pour équilibrer éthique et innovation}  
Selon Lankoski et Smith, la digitalisation étend les enjeux CR (ex. : algorithmes autonomes), les stakeholders (distants, futurs) et complique l'attribution (chaînes complexes). Des arguments des cours renforcent cela : dans les théories économiques (Williamson), les coûts de transaction (imprévisibilité) exigent une régulation organisationnelle intégrant éthique pour contrer la désorganisation, comme des sous-systèmes pour stakeholders élargis. De même, les approches empiristes (Schein) définissent la culture comme valeurs et rites, où l'intégration éthique mobilise via identité, équilibrant innovation digitale et responsabilité envers des mécanismes de clan ou bureaucratiques.

\item \textbf{Les implications pratiques pour managers, advocates, scholars et policymakers comme mécanismes d'adaptation durable}  
Le texte identifie des implications : managers intègrent CR digitale, advocates adaptent normes, scholars revisitent théories, policymakers régulent. Cela s'appuie sur le paradigme SCP (Bain), vu en cours, où la structure (digitale) influence comportements éthiques pour performance. Dans l'école des relations humaines (Bavelas), les réseaux de communication (centralité) favorisent une adaptation via interactions, suggérant que ces implications agissent comme mécanismes pour une CR durable, adaptant aux défis digitaux via apprentissage collectif et régulation.
\end{enumerate}

\subsubsection{Conclusion Persuasive}
En conclusion, la digitalisation exige une revisite de la CR via ses questions fondamentales pour une adaptation éthique, comme le démontrent Lankoski et Smith, équilibrant innovation et responsabilité. Cette vision persuasive appelle à des actions intégrées, alignées sur les théories organisationnelles des cours (contingence pour adaptation, économiques pour régulation, empiristes pour culture éthique). Ainsi, la CR digitale n'est pas un add-on mais un pilier, incitant à repenser les pratiques pour une résilience sociétale durable plutôt que purement technologique.

\section{Conclusion Générale}
Cette note de lecture combinée met en évidence des thèmes interconnectés : l'autonomie comme processus d'apprentissage, le rôle des leaders dans la culture d'apprentissage, l'innovation managériale contextuelle, la construction culturelle par les intermédiaires, et l'adaptation de la CR au digital. Ces analyses, ancrées dans les théories organisationnelles du cours, soulignent l'importance d'une approche adaptative et collective pour la performance durable dans des environnements complexes.

\end{document}